%==============================================================================
% CHAPTER 1: The Bounded Phase Space Law
%==============================================================================

\chapter{The Bounded Phase Space Law}
\label{chap:bounded_phase_space}

\epigraph{%
  The atoms of Democritus and Newton's particles of light\\
  are sands upon the Red Sea shore,\\
  where Israel's tents do shine so bright.%
}{William Blake, \textit{Mock on, Mock on, Voltaire, Rousseau} (c.~1800--1803)}

\begin{chapterbox}
This chapter derives the Bounded Phase Space Law from two axioms and establishes
its three-way equivalence theorem: oscillation, categorical distinction, and
partition operation are the same phenomenon viewed from different coordinate
systems. From the geometry of nested partitions in a bounded phase space, the
quantum numbers $(n, l, m, s)$ emerge with their exact degeneracy $C(n) = 2n^2$,
reproducing atomic shell structure without appeal to the Schr\"{o}dinger equation.
Partition depth $\Depth$ is identified with thermodynamic entropy via
$S_\Part = \kB \Depth \ln b$, and the bounds $\Depth_{\min} = \log_b 2$,
$\Depth_{\max} = \log_b(\Omega/h^n)$ establish that unobservable systems have
$\Depth = 0$. Spectroscopic selection rules and the hydrogen 21\,cm line
($\nu = 1420.405$\,MHz) are derived as partition adjacency constraints.
\end{chapterbox}

%------------------------------------------------------------------------------
\section{Two Axioms}
\label{sec:axioms}
%------------------------------------------------------------------------------

Classical statistical mechanics inherits from Liouville's theorem a picture of
phase space as an unbounded arena through which trajectories flow forever,
conserving volume but distributing themselves over arbitrarily large regions.
The ergodic hypothesis then requires the entire accessible shell of constant
energy to be explored. This picture is mathematically convenient but
physically unattainable: every laboratory system, every cell, every atom
is embedded in a finite universe and occupies a bounded region. The following
two axioms make this observation precise and draw from it consequences that
extend far beyond thermodynamics.

\begin{axiom}[Boundedness]
\label{ax:boundedness}
Every localized physical system occupies a bounded region of phase space:
\begin{equation}
  \Omega \subset \mathbb{R}^{2n}, \qquad \operatorname{Vol}(\Omega) < \infty.
\end{equation}
\end{axiom}

\begin{axiom}[Finite Partition]
\label{ax:finite_partition}
A bounded phase space $\Omega$ admits a partition into
\begin{equation}
  N = \frac{\operatorname{Vol}(\Omega)}{h^n} < \infty
\end{equation}
distinguishable cells, where $h$ is Planck's constant and the minimum cell
volume $h^n$ is set by the Heisenberg uncertainty principle.
\end{axiom}

\begin{remark}
Axiom~\ref{ax:finite_partition} is not additional structure imposed on
phase space: it follows from Axiom~\ref{ax:boundedness} together with the
fact that quantum mechanics provides a natural infrared cutoff on cell
volume. The number $N$ is the Weyl formula for the number of distinct
quantum states, recovered here from pure geometry.
\end{remark}

The finiteness of $N$ has an immediate and underappreciated consequence.

\begin{proposition}[Poincar\'{e} Recurrence is Exact]
\label{prop:poincare}
For any system satisfying Axioms~\ref{ax:boundedness}--\ref{ax:finite_partition},
Poincar\'{e} recurrence is not merely a theorem about measure-zero exceptional
trajectories but holds for \emph{every} trajectory. The recurrence time satisfies
$T_{\mathrm{rec}} \leq N \, \tau_{\min}$, where $\tau_{\min}$ is the minimum
time for a phase-space transition.
\end{proposition}

\begin{proof}
Since $N < \infty$, the phase space is a finite set of cells under
Axiom~\ref{ax:finite_partition}. Any measure-preserving map on a finite set
is a permutation, and every permutation has finite order. Therefore every
trajectory returns to its starting cell in at most $N$ steps, giving the
stated bound. The exceptional-measure caveat of the classical theorem vanishes
because there are no continuous trajectories in a discrete cell structure.
\end{proof}

%------------------------------------------------------------------------------
\section{The Triple Equivalence Theorem}
\label{sec:triple_equivalence}
%------------------------------------------------------------------------------

The central result of this chapter makes precise a conceptual unity that has
been implicit in physics since Bohr's quantization condition but never stated
as a theorem.

\begin{theorem}[Triple Equivalence]
\label{thm:triple_equivalence}
For any system satisfying Axioms~\ref{ax:boundedness}--\ref{ax:finite_partition},
the following three concepts are equivalent:
\begin{enumerate}[label=(\roman*)]
  \item \textbf{Oscillation}: the trajectory of the system in $\Omega$ is
        periodic or quasi-periodic with at least one turning point.
  \item \textbf{Categorical Distinction}: there exist at least two distinguishable
        states of the system, \ie two cells $c_1 \neq c_2$ in the partition.
  \item \textbf{Partition Operation}: the phase space $\Omega$ admits a
        non-trivial partition $\Part$ with at least two non-empty elements.
\end{enumerate}
\end{theorem}

\begin{proof}
We prove the three implications in a cycle.

\medskip
\noindent\textbf{(i) $\Rightarrow$ (iii): Oscillation implies Partition.}
Let the trajectory $\gamma(t)$ have a turning point at $q^* \in \Omega$,
meaning $\dot{q}^*(t^*) = 0$ for some $t^*$. The turning point set
$\mathcal{T} = \{q : \dot{q} = 0\} \cap \partial\Omega$ is non-empty by
hypothesis. The turning points partition $\Omega$ into regions of positive
and negative momentum: $\Omega^+ = \{(q,p) : p > 0\}$,
$\Omega^- = \{(q,p) : p < 0\}$, and $\mathcal{T}$ itself. Since the trajectory
must reverse direction at each turning point (by continuity of the Hamiltonian
flow and $\operatorname{Vol}(\Omega) < \infty$), the set
$\Part = \{\Omega^+, \Omega^-, \mathcal{T}\}$ is a non-trivial partition of $\Omega$.

\medskip
\noindent\textbf{(iii) $\Rightarrow$ (ii): Partition implies Distinction.}
Let $\Part = \{c_1, \ldots, c_N\}$ with $N \geq 2$. By Axiom~\ref{ax:finite_partition},
distinct cells have non-overlapping support and both have positive measure. A
physical observable $O$ that takes different values on $c_1$ and $c_2$ (for
example, the cell index itself, or equivalently the number of completed
partition operations from any reference state) witnesses that $c_1$ and $c_2$
are categorically distinct. The existence of such an observable is guaranteed
by the separability of $\mathbb{R}^{2n}$.

\medskip
\noindent\textbf{(ii) $\Rightarrow$ (i): Distinction implies Oscillation.}
Suppose cells $c_1 \neq c_2$ are distinguishable states. In a bounded phase
space, the system trajectory must remain within $\Omega$. If the trajectory
reaches $c_2$ starting from $c_1$, it cannot escape to infinity; it must
return toward the interior of $\Omega$. By Proposition~\ref{prop:poincare},
every trajectory in a finite cell structure recurs. Recurrence of a
trajectory that passes through both $c_1$ and $c_2$ requires it to reverse
the transition $c_1 \to c_2$, which constitutes a turning point and hence
oscillation. \qed
\end{proof}

\begin{remark}
The proof of (ii)~$\Rightarrow$~(i) uses recurrence essentially: in an unbounded
phase space, a system could distinguish two states and escape to infinity,
never returning. Boundedness closes this escape route, forging the equivalence.
This is why free particles (formally unbounded) do not oscillate in the absence
of a confining potential.
\end{remark}

The connection to Birkhoff's ergodic theorem is instructive. Birkhoff establishes
that time averages equal space averages for almost all initial conditions on
ergodic components. In a finite partition, every ergodic component is itself
a finite collection of cells, so the Birkhoff average is a finite sum. The
partition depth $\Depth$ (defined in Section~\ref{sec:partition_depth}) serves
as a Kolmogorov--Sinai entropy rate: $h_{\mathrm{KS}} = \Depth \ln b / \tau_{\min}$,
connecting the abstract combinatorics of partitions to measurable dynamical
quantities.

%------------------------------------------------------------------------------
\section{Partition Coordinates and Shell Structure}
\label{sec:partition_coordinates}
%------------------------------------------------------------------------------

A bounded phase space $\Omega$ admits not merely one partition but an entire
hierarchy of nested partitions. This hierarchy is the central object of
categorical mechanics.

\begin{definition}[Nested Partition Hierarchy]
\label{def:nested_partition}
A nested partition hierarchy on $\Omega$ is a sequence
\[
  \Omega = \Part_0 \supset \Part_1 \supset \Part_2 \supset \cdots
\]
where each $\Part_k$ is a refinement of $\Part_{k-1}$, meaning every element
of $\Part_k$ is contained in some element of $\Part_{k-1}$.
\end{definition}

Each element of $\Part_k$ is characterized by four integer-valued coordinates
that track its position in the hierarchy.

\begin{definition}[Partition Coordinates]
\label{def:partition_coordinates}
A state in the nested partition hierarchy of $\Omega$ is characterized by the
quadruple $(n, l, m, s)$ where:
\begin{itemize}
  \item $n \geq 1$: \emph{nesting depth}, the level in the hierarchy at which
        this state first becomes a distinct element.
  \item $l \in \{0, 1, \ldots, n-1\}$: \emph{partition complexity}, the number
        of non-trivial sub-partitions within the current element.
  \item $m \in \{-l, -l+1, \ldots, +l\}$: \emph{partition orientation},
        the signed index of the sub-partition within its complexity class.
  \item $s \in \{-\tfrac{1}{2}, +\tfrac{1}{2}\}$: \emph{intrinsic parity},
        the orientation of the element with respect to its parent partition
        boundary (interior vs.\ boundary-adjacent).
\end{itemize}
\end{definition}

\begin{theorem}[Shell Capacity from Partition Geometry]
\label{thm:shell_capacity}
The number of distinct states at nesting depth $n$ in the partition hierarchy
of a bounded phase space is:
\begin{equation}
  C(n) = 2n^2.
\end{equation}
The sequence $C(1), C(2), C(3), C(4), \ldots = 2, 8, 18, 32, \ldots$ is
the exact sequence of atomic shell capacities.
\end{theorem}

\begin{proof}
At depth $n$, the complexity $l$ ranges over $\{0, 1, \ldots, n-1\}$,
giving $n$ choices. For each $l$, the orientation $m$ ranges over the
$2l+1$ values $\{-l, \ldots, +l\}$. For each $(l, m)$ pair, the parity
$s$ contributes a factor of 2. Therefore:
\begin{align}
  C(n) &= 2 \sum_{l=0}^{n-1} (2l+1) \notag \\
       &= 2 \sum_{l=0}^{n-1} (2l+1) \notag \\
       &= 2 \Bigl[ 2\cdot\frac{(n-1)n}{2} + n \Bigr] \notag \\
       &= 2 \bigl[ n(n-1) + n \bigr] \notag \\
       &= 2n^2. \qquad\blacksquare
\end{align}
\end{proof}

\begin{remark}[On the Non-Status of Schr\"{o}dinger's Equation Here]
Theorem~\ref{thm:shell_capacity} derives $C(n) = 2n^2$ from the combinatorics
of nested partitions in a bounded phase space. The Schr\"{o}dinger equation
is not invoked. The quantum numbers $(n, l, m, s)$ are labels for positions
in a partition hierarchy; their range constraints and degeneracies follow from
the geometry of nesting, not from boundary conditions on a differential equation.
That the Schr\"{o}dinger equation yields the same degeneracies is not a
coincidence: it is a different coordinate system for the same underlying
partition structure.
\end{remark}

%------------------------------------------------------------------------------
\section{Partition Depth and Its Equivalence to Entropy}
\label{sec:partition_depth}
%------------------------------------------------------------------------------

\begin{definition}[Partition Depth]
\label{def:partition_depth}
Let $\Part$ be a partition generated by a sequential process with branching
factors $k_1, k_2, \ldots$ at successive levels and encoding base $b \geq 2$.
The \emph{partition depth} is:
\begin{equation}
  \Depth = \sum_i \log_b k_i.
  \label{eq:partition_depth}
\end{equation}
\end{definition}

The partition depth $\Depth$ measures the total amount of categorical information
generated by all partitioning operations. Each level with branching factor $k_i$
contributes $\log_b k_i$ units of depth, which is the information content of
a single choice among $k_i$ alternatives encoded in base $b$.

\begin{theorem}[Depth--Entropy Equivalence]
\label{thm:depth_entropy}
The thermodynamic entropy of a partition $\Part$ with depth $\Depth$ is:
\begin{equation}
  S_\Part = \kB \Depth \ln b.
  \label{eq:depth_entropy}
\end{equation}
\end{theorem}

\begin{proof}
Consider a partition $\Part$ with $W$ equally likely cells, achieved through
a sequence of branching operations with factors $k_1, \ldots, k_r$ such that
$W = \prod_i k_i$. By Boltzmann's relation, $S = \kB \ln W$.
Now $\ln W = \ln \prod_i k_i = \sum_i \ln k_i$.
Converting to base $b$: $\sum_i \ln k_i = \ln b \sum_i \log_b k_i = \ln b \cdot \Depth$
by Definition~\ref{def:partition_depth}. Substituting:
\begin{equation}
  S_\Part = \kB \ln W = \kB \ln b \cdot \Depth = \kB \Depth \ln b. \qquad\blacksquare
\end{equation}
\end{proof}

\begin{remark}[Identification with Shannon Information]
For $b = 2$, Equation~\eqref{eq:depth_entropy} gives $S_\Part = \kB \ln 2 \cdot \Depth$.
Since $\Depth$ is measured in bits (log base 2), and Shannon entropy is
$H = \Depth$ bits, we have $S_\Part = \kB \ln 2 \cdot H$. This is precisely
the Brillouin relation between physical entropy and information content,
here derived from first principles rather than postulated.
\end{remark}

The depth is bounded above and below.

\begin{theorem}[Depth Bounds]
\label{thm:depth_bounds}
For any system satisfying Axioms~\ref{ax:boundedness}--\ref{ax:finite_partition}:
\begin{equation}
  \Depth_{\min} = \log_b 2 \leq \Depth \leq \log_b\!\left(\frac{\operatorname{Vol}(\Omega)}{h^n}\right) = \Depth_{\max}.
  \label{eq:depth_bounds}
\end{equation}
\end{theorem}

\begin{proof}
\textbf{Lower bound.} By the Triple Equivalence Theorem~\ref{thm:triple_equivalence},
a system with $\Depth < \log_b 2$ would have fewer than $b^{\log_b 2} = 2$
distinguishable states, meaning no partition exists with more than one cell.
Such a system is categorically trivial---it has no internal distinctions,
hence no partition, hence (by the equivalence) no oscillation and no observable
dynamics. The minimum non-trivial partition has exactly 2 cells (one binary
partition operation), giving $\Depth_{\min} = \log_b 2$.

\textbf{Upper bound.} By Axiom~\ref{ax:finite_partition}, the total number
of distinguishable cells is $N = \operatorname{Vol}(\Omega)/h^n < \infty$.
The maximum partition depth occurs when every cell is individually distinguished,
which requires $\Depth_{\max} = \log_b N = \log_b(\operatorname{Vol}(\Omega)/h^n)$.
No further refinement is physically meaningful because sub-Planck cells
are indistinguishable by the uncertainty principle. \qed
\end{proof}

\begin{corollary}[Existence Requires Depth]
\label{cor:existence_depth}
A system with $\Depth = 0$ is unobservable. Physical existence as a distinct
entity requires $\Depth \geq \Depth_{\min} = \log_b 2$.
\end{corollary}

\begin{proof}
$\Depth = 0$ implies $W = b^0 = 1$: exactly one state. A system with one state
cannot be distinguished from the background; it makes no observable difference
to any measurement apparatus (which would itself require $\Depth \geq \log_b 2$
to register a result). Hence $\Depth = 0$ corresponds to physical
non-existence as a distinct entity. The first detectable level is $\Depth = \log_b 2$,
the minimum single partition operation. \qed
\end{proof}

The existence cost $\Eexist$ introduced in later chapters is the energetic
price of maintaining $\Depth \geq \log_b 2$ against dissipation:
$\Eexist = \kB T \ln 2$ per partition boundary, recoverable as Landauer's bound
in the special case of a one-bit distinction.

%------------------------------------------------------------------------------
\section{Spectroscopic Selection Rules as Partition Adjacency}
\label{sec:selection_rules}
%------------------------------------------------------------------------------

The partition coordinates $(n, l, m, s)$ of Definition~\ref{def:partition_coordinates}
carry their own geometry: two states are \emph{adjacent} in the partition
hierarchy if and only if they differ by a single elementary partition operation.

\begin{definition}[Partition Adjacency]
\label{def:partition_adjacency}
Two states $(n, l, m, s)$ and $(n', l', m', s')$ are \emph{partition-adjacent}
if the partition operation connecting them changes exactly one coordinate by
the minimum possible amount:
\begin{equation}
  |n' - n| \leq 1, \quad |l' - l| = 1, \quad |m' - m| \leq 1, \quad s' = s.
\end{equation}
\end{definition}

\begin{theorem}[Selection Rules from Adjacency]
\label{thm:selection_rules}
The spectroscopic electric dipole selection rules
\begin{equation}
  \Delta l = \pm 1, \qquad \Delta m \in \{0, \pm 1\}, \qquad \Delta s = 0
  \label{eq:selection_rules}
\end{equation}
are exactly the conditions for partition adjacency in the hierarchy of
Definition~\ref{def:nested_partition}.
\end{theorem}

\begin{proof}
An electric dipole photon carries angular momentum quantum numbers
$(j=1, m_j \in \{-1, 0, +1\})$. In the partition framework, photon emission
or absorption corresponds to a single partition operation: one distinction is
either created ($\Delta l = +1$) or merged ($\Delta l = -1$). A single
partition operation changes the complexity $l$ by exactly $\pm 1$.

The orientation $m$ must change by the minimum required to accommodate the
partition operation within the new complexity class. Since the photon contributes
$\Delta m_j \in \{-1, 0, +1\}$ and total orientation is conserved modulo the
partition operation, $\Delta m \in \{-1, 0, +1\}$ follows.

The parity $s$ is a property of the parent partition structure, not of the
individual partition operation, so it is unchanged: $\Delta s = 0$.

The forbidden transition $\Delta l = 0$ for dipole radiation corresponds to
a ``null partition operation''---a change that creates no new distinction---which
is not a partition operation at all. The Laporte rule ($\Delta l = \pm 1$ with
parity change for centrosymmetric systems) is the statement that a single
partition operation strictly changes the complexity. \qed
\end{proof}

\begin{example}[Hydrogen 21\,cm Line]
\label{ex:21cm}
The hydrogen hyperfine transition at $\nu = 1420.405\,\mathrm{MHz}$
(wavelength $\lambda \approx 21\,\mathrm{cm}$, energy splitting
$\Delta E = h\nu \approx 5.87 \times 10^{-6}\,\mathrm{eV}$) arises from
the flip of the electron spin parity $s$ relative to the proton spin parity
in the ground state $n = 1$, $l = 0$, $m = 0$.

In the partition framework, the hydrogen ground state has $\Depth_{\min} = \log_2 2 = 1$
with a two-cell partition corresponding to the two parity states
$(s = +\tfrac{1}{2}, s = -\tfrac{1}{2})$. The hyperfine coupling
$A_{\mathrm{HF}} = h\nu_{\mathrm{HF}}$ is the energy cost of a parity-flip
operation---the minimum partition operation on the intrinsic-parity coordinate.

The transition $\Delta s = \pm 1$ (triplet $\to$ singlet) with $\Delta l = 0$,
$\Delta m = 0$ is a magnetic dipole transition, consistent with
Definition~\ref{def:partition_adjacency} because here it is $s$ (parity),
not $l$ (complexity), that changes by its minimum step. This is the
complementary adjacency to the electric dipole case: electric dipole changes
$l$; magnetic dipole changes $s$. The numerical value of $\nu = 1420.405\,\mathrm{MHz}$
encodes the energy scale of the proton--electron partition coupling in the
ground state, and its extraordinary precision (the most precisely known
radio frequency in astronomy) reflects the rigidity of the minimum-depth
partition structure.
\end{example}

%------------------------------------------------------------------------------
\section{Connection to the Birkhoff--Kolmogorov Framework}
\label{sec:birkhoff_kolmogorov}
%------------------------------------------------------------------------------

The partition depth $\Depth$ is related to the Kolmogorov--Sinai (KS) entropy
of the dynamical system in a precise way.

\begin{proposition}[Depth as KS Entropy Rate]
\label{prop:ks_entropy}
For a dynamical system on a bounded phase space $\Omega$ with minimum transition
time $\tau_{\min}$, the Kolmogorov--Sinai entropy satisfies:
\begin{equation}
  h_{\mathrm{KS}} = \frac{\Depth \ln b}{\tau_{\min}}.
  \label{eq:ks_entropy_rate}
\end{equation}
\end{proposition}

\begin{proof}
The KS entropy of a partition $\Part$ is defined as
$h_{\mathrm{KS}} = \lim_{T\to\infty} \frac{1}{T} H(\Part^T)$
where $\Part^T$ is the join of all time-translates of $\Part$ up to time $T$,
and $H$ is the Shannon entropy in nats. By Theorem~\ref{thm:depth_entropy},
$H(\Part) = \Depth \ln b$ for the generating partition. In a bounded system
with $N < \infty$ cells and minimum transition time $\tau_{\min}$, the number
of distinct $T$-length orbits saturates at $N$ for $T \gg N\tau_{\min}$.
The entropy rate therefore equals the per-step entropy divided by $\tau_{\min}$:
$h_{\mathrm{KS}} = \Depth \ln b / \tau_{\min}$. \qed
\end{proof}

Birkhoff's ergodic theorem guarantees that for ergodic systems, the time
average of any observable equals its space average. In the partition framework,
this means that the time average of the partition depth $\Depth(t)$ equals
the microcanonical average over all cells at the appropriate energy shell.
This is not merely an abstract result: it provides the foundation for
Chapter~\ref{chap:charge_emergence}, where the time-averaged partition depth
of a bound state is identified with the bond energy.

\begin{proposition}[Birkhoff Average of Depth]
\label{prop:birkhoff_depth}
For an ergodic system on a bounded phase space, the time-averaged partition
depth satisfies:
\begin{equation}
  \langle \Depth \rangle_t = \frac{1}{N} \sum_{c \in \Part} \Depth(c),
  \label{eq:birkhoff_depth}
\end{equation}
where the sum is over all cells $c$ of the generating partition $\Part$.
\end{proposition}

The right-hand side of \eqref{eq:birkhoff_depth} is the categorical richness
$\Rich = N^{-1} \sum_c \Depth(c)$, a quantity that will play a central role
in describing the information content of biological sequences in Part~IV.

%------------------------------------------------------------------------------
\section{Discussion: What This Replaces and Why}
\label{sec:discussion_ch1}
%------------------------------------------------------------------------------

The framework developed in this chapter reproduces, from two axioms, the
following results that are conventionally taken as empirical inputs or as
outputs of the Schr\"{o}dinger equation:

\begin{enumerate}[label=(\arabic*)]
  \item \textbf{Quantum numbers and their ranges}: $(n, l, m, s)$ emerge as
        coordinates in a nested partition hierarchy, with their exact domains
        $l \in \{0, \ldots, n-1\}$, $m \in \{-l, \ldots, +l\}$,
        $s \in \{-\tfrac{1}{2}, +\tfrac{1}{2}\}$ following from the structure
        of the hierarchy (Theorem~\ref{thm:shell_capacity}).

  \item \textbf{Atomic shell capacities}: $C(n) = 2n^2$ is a partition counting
        identity (Theorem~\ref{thm:shell_capacity}), not a consequence of the
        Pauli exclusion principle, though the exclusion principle is consistent
        with it: no two electrons can occupy the same partition cell.

  \item \textbf{Entropy--information equivalence}: $S = \kB \ln 2 \cdot H$
        (Brillouin relation) is a special case of Theorem~\ref{thm:depth_entropy},
        not a separate postulate.

  \item \textbf{Spectroscopic selection rules}: $\Delta l = \pm 1$,
        $\Delta m \in \{0, \pm 1\}$ are partition adjacency conditions
        (Theorem~\ref{thm:selection_rules}), making them topological rather
        than dynamical constraints.

  \item \textbf{Landauer's bound}: $E_{\mathrm{exist}} = \kB T \ln 2$ per bit
        is the special case of the depth--entropy relation at $\Depth = \log_2 2 = 1$,
        confirming that the minimum energetic cost of a physical distinction is
        the Boltzmann unit of information.
\end{enumerate}

What this framework does \emph{not} replace is quantitative prediction of energy
levels: the precise value of $\Delta E$ for a given transition requires the
Schr\"{o}dinger equation (or its relativistic generalization) applied to a
specific Hamiltonian. The partition framework provides the topology---which
transitions exist and which are forbidden---while the Hamiltonian provides
the metric, assigning energy to each partition cell. The two are complementary,
not competing.

The 21\,cm line (Example~\ref{ex:21cm}) illustrates this division: the
\emph{existence} and \emph{character} of the transition (magnetic dipole,
minimum parity flip) follow from partition geometry; the precise frequency
$1420.405\,\mathrm{MHz}$ requires the Fermi contact interaction and QED
corrections. Partition theory predicts the graph; quantum field theory fills
in the edge weights.

Chapter~\ref{chap:charge_emergence} extends this framework to ask: what
is the \emph{boundary} of a partition? The answer---that partition boundaries
carry the signature of electric charge---is the second major result of
Part~\ref{part:foundations}.

%------------------------------------------------------------------------------
\section{Categorical Distance and the Geometry of State Space}
\label{sec:categorical_distance}
%------------------------------------------------------------------------------

The partition coordinates of Section~\ref{sec:partition_coordinates} do more
than label states; they endow the state space with a natural metric.

\begin{definition}[Categorical Distance]
\label{def:categorical_distance}
For two states $\sigma_1 = (n_1, l_1, m_1, s_1)$ and
$\sigma_2 = (n_2, l_2, m_2, s_2)$ in the partition hierarchy, the
\emph{categorical distance} is:
\begin{equation}
  \dC(\sigma_1, \sigma_2) = |n_1 - n_2| + |l_1 - l_2| + |m_1 - m_2|
    + |s_1 - s_2|,
  \label{eq:categorical_distance}
\end{equation}
where the chirality difference $|s_1 - s_2| \in \{0, 1\}$.
\end{definition}

\begin{proposition}[Categorical Distance is a Metric]
\label{prop:metric}
The categorical distance $\dC$ satisfies the axioms of a metric on the set
of partition coordinates:
\begin{enumerate}[label=(\alph*)]
  \item \textbf{Non-negativity}: $\dC(\sigma_1, \sigma_2) \geq 0$, with
        equality iff $\sigma_1 = \sigma_2$.
  \item \textbf{Symmetry}: $\dC(\sigma_1, \sigma_2) = \dC(\sigma_2, \sigma_1)$.
  \item \textbf{Triangle inequality}: $\dC(\sigma_1, \sigma_3) \leq
        \dC(\sigma_1, \sigma_2) + \dC(\sigma_2, \sigma_3)$.
\end{enumerate}
\end{proposition}

\begin{proof}
(a) Each absolute value term is non-negative. Equality to zero requires all
four terms to vanish simultaneously, which is equivalent to $\sigma_1 = \sigma_2$.
(b) Immediate from the symmetry of absolute value.
(c) The triangle inequality for the integer $L^1$ (Manhattan) distance is a
standard result: $|n_1 - n_3| \leq |n_1 - n_2| + |n_2 - n_3|$, and similarly
for each coordinate. Summing over all four coordinates gives (c). \qed
\end{proof}

\begin{corollary}[Selection Rules as Unit-Ball Constraint]
\label{cor:unit_ball}
The electric dipole selection rules (Theorem~\ref{thm:selection_rules}) are
exactly the condition that a transition is allowed only between states with
$\dC = 1$ in the $(l, m, s)$ subspace:
\begin{equation}
  \dC^{(l,m,s)}(\sigma_1, \sigma_2) = 1 \iff \text{electric dipole transition allowed.}
\end{equation}
\end{corollary}

\begin{proof}
From Theorem~\ref{thm:selection_rules}: $\Delta l = \pm 1$ contributes
$|l_1 - l_2| = 1$; $\Delta m \in \{0, \pm 1\}$ contributes $|m_1 - m_2| \leq 1$;
$\Delta s = 0$ contributes $0$. The total distance in the $(l, m)$ subspace
is $1 + |m_1 - m_2|$, which equals $1$ only when $\Delta m = 0$.
For $\Delta m = \pm 1$, the distance is $2$. However, the $\Delta l$ change
is mandatory (contributes exactly 1), so the relevant condition is $|\Delta l| = 1$,
which corresponds to $\dC^{(l)}(\sigma_1,\sigma_2) = 1$. The allowed $\Delta m$
values $\{0, \pm 1\}$ are the unit ball condition on the $m$-subspace given
that $\Delta l = 1$ must hold. \qed
\end{proof}

The categorical distance measures the ``cost'' of a transition in partition
space. Transitions with large $\dC$ require multiple partition operations and
are correspondingly suppressed (higher-multipole radiation, Raman processes).
This gives the partition framework's account of why selection rules are selection
\emph{rules} and not merely selection \emph{tendencies}: a single photon
interacts with a single partition boundary, so it can only change one
partition coordinate by its minimum step.

%------------------------------------------------------------------------------
\section{Existence Cost and Landauer's Bound}
\label{sec:existence_cost}
%------------------------------------------------------------------------------

Every categorical distinction carries an energetic price. This is the content
of Landauer's principle, which in the partition framework is not a separate
postulate but a corollary of the depth--entropy equivalence.

\begin{theorem}[Existence Cost]
\label{thm:existence_cost}
The minimum energy required to maintain one categorical distinction---one
non-trivial partition boundary---at temperature $T$ is:
\begin{equation}
  \Eexist = \kB T \ln 2.
  \label{eq:existence_cost}
\end{equation}
This is Landauer's bound, derived from the Bounded Phase Space Axioms.
\end{theorem}

\begin{proof}
By Corollary~\ref{cor:existence_depth}, the minimum partition depth for a
physically existent entity is $\Depth_{\min} = \log_b 2$, corresponding to
a single binary partition. By Theorem~\ref{thm:depth_entropy}:
\begin{equation}
  S_{\min} = \kB \Depth_{\min} \ln b = \kB \log_b 2 \cdot \ln b = \kB \ln 2.
\end{equation}
At temperature $T$, the free energy cost of maintaining this entropy state
against thermal equilibration (which tends to $\Depth = 0$ under irreversible
diffusion) is $F = -TS_{\min} = -T \kB \ln 2$; the energy \emph{expended}
to maintain the distinction is therefore $\Eexist = \kB T \ln 2$ per erasure
cycle. This is exactly Landauer's bound. \qed
\end{proof}

\begin{remark}[Physical Interpretation]
At biological temperature $T = 310\,\mathrm{K}$:
$\Eexist = 1.381\times10^{-23} \times 310 \times \ln 2 \approx 2.97\times10^{-21}\,\mathrm{J}$
$\approx 0.018\,\mathrm{eV}$. This is significantly smaller than a typical
covalent bond energy ($\sim 3$--$5\,\mathrm{eV}$) but of the same order as
hydrogen bond energies ($\sim 0.1$--$0.3\,\mathrm{eV}$). The existence cost
$\Eexist$ sets the minimum energy scale for a biological information bit---a
result that will be important when analyzing the energetics of genomic computation
in Part~IV.
\end{remark}

\begin{corollary}[Bennett's Reversible Computing]
\label{cor:bennett}
A partition operation is logically reversible if and only if it preserves
partition depth $\Depth$, which is if and only if it maps one partition to
another of equal cardinality. In this case, the thermodynamic cost
\eqref{eq:existence_cost} is zero; irreversibility is the cost of depth
\emph{erasure}, not of the logical operation itself.
\end{corollary}

\begin{proof}
Reversible computation preserves information, hence preserves $\Depth$.
By Theorem~\ref{thm:depth_entropy}, $\Delta S = 0$ for a reversible computation,
so $\Delta F = 0$: no energy need be dissipated. The Landauer cost arises only
when a partition is \emph{merged} (depth reduced), which is irreversible.
This is Bennett's result on reversible computation, derived here from the
partition axioms. \qed
\end{proof}

%=============================================================================
\section*{Chapter Summary}
\label{sec:ch1_summary}
%=============================================================================

\begin{chapterbox}[title=Chapter 1 — Key Results Established]
\begin{enumerate}[label=\textbf{\arabic*.}]
  \item \textbf{Two Axioms}: Boundedness ($\operatorname{Vol}(\Omega) < \infty$)
        and Finite Partition ($N = \operatorname{Vol}(\Omega)/h^n < \infty$)
        are sufficient for all results in this chapter. No dynamical equations
        are assumed.

  \item \textbf{Triple Equivalence Theorem~\ref{thm:triple_equivalence}}:
        Oscillation $\equiv$ Categorical Distinction $\equiv$ Partition Operation.
        In bounded phase space these three concepts are not merely related---they
        are identical descriptions of the same phenomenon.

  \item \textbf{Poincar\'{e} Recurrence Theorem~\ref{prop:poincare}}:
        In a finite partition, recurrence is exact for every trajectory,
        not just measure-almost-all. Recurrence time $T_{\mathrm{rec}} \leq N\tau_{\min}$.

  \item \textbf{Shell Capacity Theorem~\ref{thm:shell_capacity}}:
        $C(n) = 2n^2$, derived from partition combinatorics. The sequence
        $2, 8, 18, 32, \ldots$ is a geometric necessity, not an empirical fact.

  \item \textbf{Depth--Entropy Equivalence Theorem~\ref{thm:depth_entropy}}:
        $S_\Part = \kB \Depth \ln b$. Partition depth is thermodynamic
        entropy in a different coordinate.

  \item \textbf{Depth Bounds Theorem~\ref{thm:depth_bounds}}:
        $\log_b 2 \leq \Depth \leq \log_b(\operatorname{Vol}(\Omega)/h^n)$.
        Unobservable systems have $\Depth = 0$.

  \item \textbf{Selection Rules Theorem~\ref{thm:selection_rules}}:
        $\Delta l = \pm 1$, $\Delta m \in \{0,\pm 1\}$, $\Delta s = 0$
        are partition adjacency conditions, not dynamical constraints.

  \item \textbf{KS Entropy Rate Proposition~\ref{prop:ks_entropy}}:
        $h_{\mathrm{KS}} = \Depth \ln b / \tau_{\min}$, connecting the
        combinatorial depth to measurable dynamical entropy production.

  \item \textbf{Existence Cost Theorem~\ref{thm:existence_cost}}:
        $\Eexist = \kB T \ln 2$ is Landauer's bound, derived from the axioms
        as the minimum energy for one categorical distinction.

  \item \textbf{Complementarity}: The partition framework provides the
        \emph{topology} of physical law---which states exist, which transitions
        are allowed, what selection rules hold---while dynamical equations
        (Schr\"{o}dinger, Maxwell) provide the \emph{metric}, assigning
        energies and timescales to the topological skeleton.
\end{enumerate}
\end{chapterbox}
